% Evidence appendix, part 1 of 3 (continued in parts 2 and 3); label prefix: evi. No new macros.
\section{Evidence: finite computations}
\label{sec:evi}
\noindent\genusstatus{numerically supported} The finite computations recorded in this appendix are exact or modular computations; they support the statements at the genera, orders and primes given, and are not proofs. The code and outputs are in \texttt{code/} of this project.


Nothing in this appendix is a step of a proof: the proofs of the note are complete without it, and no proof refers to it. It records in tables some finite computations made while the statements were developed, with exact arithmetic over $\bbQ$ or $\bbQ(\zeta)$ or arithmetic modulo the primes named (a zero test modulo a prime is probabilistic); conifold frames are labelled by $(\alpha,\beta)$ as in \cref{def:set-frames}, $\varpi_0=f_0+\alpha\varpi_1+\beta f_2$.

\subsection{Graph sums of genus 2 to 5}
\label{sec:evi-graphs}

The outer terms $T_\Gamma=u^{2-2g}G_\Gamma(u)$ (\cref{def:set-outer,lem:flh-hat}), with $M=\Mo(\tau(u)\zh)$, were computed for every leg-free stable graph of genus $2$ to $5$, modulo $2^{89}-1$ and $2^{61}-1$ with rational reconstruction, in the frames $(0,0)$ and $(1/3,-2/7)$; an edge graph is regular if and only if $G_\Gamma=O(u^{2g-2})$. At genus $4$ and $5$ the outer sum was obtained by regrouping the plain graph sum along the edges of $M$, a regrouping checked against genuine vertex tensors only at genus $2$ and $3$. With the bare symbol $\Mo(\zh)$ all $149$ edge graphs of genus $4$ with at most five edges are singular, and with $\tau$ replaced by $\tau+u^5$ exactly the three one-edge graphs of genus $4$ fail. The value $c_g$ and the gap of the total do not depend on $M$; only the regularity of the edge graphs tests the factorisation. The results are in \cref{tab:evi-g2,tab:evi-g3,tab:evi-g4,tab:evi-g5}.

\begin{table}[ht]
\centering\small\setlength{\tabcolsep}{4pt}
\begin{tabular}{l|r|r|r}
graph of genus $2$ & plain, $\coef{\tc^{-2}}$ & plain, $\coef{\tc^{-1}}$ & natural, at $u=0$\\ \hline
one vertex of genus $2$ & $1807/3072000$ & $-791/15360000$ & $\tfrac{1}{1200}\tc^{-2}+O(1)$\\
genus-$1$ vertex with a loop & $-811/256000$ & $-1771/3840000$ & $479/756000$\\
genus-$0$ vertex with two loops & $-57/64000$ & $41/12800$ & $1/2250$\\
two genus-$1$ vertices, one edge & $843/1024000$ & $743/15360000$ & $-127/2016000$\\
genus $0$ with a loop, edge to genus $1$ & $3/51200$ & $-41/153600$ & $0$\\
two genus-$0$ vertices, three edges & $87/32000$ & $37/160000$ & $0$\\
two genus-$0$ vertices with loops, an edge & $9/12800$ & $-173/64000$ & $0$\\ \hline
sum & $1/1200$ & $0$ & \\
\end{tabular}
\caption{Genus $2$, per stable graph, summed over colourings, in the frame $(0,0)$, modulo both primes: the polar coefficients of the plain graph sum of \cref{def:set-graphsum} (vertices the TFT $\TFT$), and the terms $T_\Gamma$ of the natural form at $u=0$ (in the frame $(1/3,-2/7)$, modulo $2^{89}-1$, the edge graphs are regular as well). The $115$ coloured graphs give identical tables at both primes.}
\label{tab:evi-g2}
\end{table}

\begin{table}[ht]
\centering\small
\begin{tabular}{l|p{0.68\textwidth}}
stable graphs & $42$; with $0,1,\dots,6$ edges: $1,2,5,9,12,8,5$\\
edge graphs & all $41$ regular: frame $(0,0)$ modulo both primes, frame $(1/3,-2/7)$ modulo $2^{89}-1$\\
vertex graph, $(0,0)$ & $u^{-4}\bigl(\tfrac{1}{25200}-\tfrac{u}{9000}+\tfrac{407u^2}{3780000}-\tfrac{151u^3}{3780000}+\dots\bigr)=\tfrac{1}{25200}\tc^{-4}+O(1)$\\
vertex graph, $(1/3,-2/7)$ & $u^{-4}\bigl(\tfrac{1}{25200}-\tfrac{11u}{189000}+\dots\bigr)$, with the same polar part\\
\end{tabular}
\caption{Genus $3$, the natural form.}
\label{tab:evi-g3}
\end{table}

\begin{table}[ht]
\centering\small
\begin{tabular}{l|p{0.72\textwidth}}
stable graphs & $379$; with $0,1,\dots,9$ edges: $1,3,7,21,43,75,89,81,42,17$\\
edge graphs & all $378$ regular, in both frames, modulo both primes\\
total, $(0,0)$ & $u^{-6}\bigl(\tfrac{1}{180000}-\tfrac{7u}{300000}+\tfrac{701u^2}{18000000}-\tfrac{1459u^3}{45000000}+\tfrac{923639u^4}{67500000000}-\tfrac{171583u^5}{67500000000}\bigr)$ $=\tfrac{1}{180000}\tc^{-6}+O(1)$\\
total, $(1/3,-2/7)$ & $u^{-6}\bigl(\tfrac{1}{180000}-\tfrac{11u}{900000}-\tfrac{79u^2}{378000000}+\dots\bigr)$, with the same polar part\\
\end{tabular}
\caption{Genus $4$, the natural form; $1/180000=c_4$.}
\label{tab:evi-g4}
\end{table}

\begin{table}[ht]
\centering\small
\begin{tabular}{l|p{0.72\textwidth}}
stable graphs & $4555$; with $0,1,\dots,12$ edges: $1,3,11,34,100,239,492,784,1002,926,632,260,71$\\
edge graphs & all $4554$ regular: frame $(0,0)$ modulo both primes, frame $(1/3,-2/7)$ modulo $2^{89}-1$\\
total, $(0,0)$ & $u^{-8}\bigl(\tfrac{1}{660000}-\tfrac{7u}{825000}+\tfrac{199u^2}{9900000}-\tfrac{717u^3}{27500000}+\tfrac{449357u^4}{22500000000}-\tfrac{5593217u^5}{618750000000}$ $+\tfrac{14686042363u^6}{6496875000000000}-\tfrac{8370858511u^7}{32484375000000000}\bigr)=\tfrac{1}{660000}\tc^{-8}+O(1)$, by the Chinese remainder theorem from both primes\\
total, $(1/3,-2/7)$ & $u^{-8}\bigl(\tfrac{1}{660000}-\tfrac{u}{225000}+\tfrac{323u^2}{207900000}+\tfrac{66841u^3}{15592500000}-\dots\bigr)$, modulo both primes: different from the first at every order $u^1,\dots,u^7$, with the same polar part\\
\end{tabular}
\caption{Genus $5$, the natural form; $1/660000=c_5$.}
\label{tab:evi-g5}
\end{table}

% Evidence appendix, part 2 of 3: the normalisation. Continued in part 3.

\subsection{The normalisation and the identity \texorpdfstring{$(\Gamma)$}{(Gamma)}}
\label{sec:evi-normal}

\Cref{tab:evi-windows} lists finite windows in which $\Dn=1$, the identity \eqref{eq:gam-Gamma} and some identities entering their proofs were tested directly, independently of the proofs of \cref{thm:gam-D} and \cref{thm:gam-main}. Since $\Phican=\PhiSP$ in every conifold frame (\cref{thm:con-Pb}) and $\PhiSP$ involves no frame, further conifold frames give no independent evidence for $\Dn=1$.

\begin{table}[ht]
\centering\small
\begin{tabular}{>{\raggedright\arraybackslash}p{0.4\textwidth}|>{\raggedright\arraybackslash}p{0.3\textwidth}|>{\raggedright\arraybackslash}p{0.22\textwidth}}
statement & range & arithmetic\\ \hline
$\Dn_0=1$, order by order at the conifold point & $z^{204}$ ($g\leq 69$), two implementations; $z^{125}$, a third & exact, frame $(0,0)$\\
$\Dn_0=1$ & $z^{1004}$ ($g\leq 335$) & modulo $2^{89}-1$, $2^{61}-1$\\
the actual gauge is $\Pc$ (\cref{cor:mai-actual}) & $z^{35}$; $z^{150}$ & two primes, two frames\\
\eqref{eq:gam-Gamma}: $c_n(1)=h_n$ & $n\leq 60$ (Laurent expansion at $L=1$); $n\leq 40$ (tracked precision); $n\leq 30$ (fit at $32$ points) & exact over $\bbQ$\\
\eqref{eq:gam-Gamma} & $z^{500}$ & modulo $2^{127}-1$, $2^{107}-1$\\
the chart of \cref{thm:gam-lmi} at $L=1$ equals $h$ & $z^{60}$ & exact over $\bbQ(\zeta)$\\
$\PhiSP=\Phican$ (\cref{thm:con-Pb}) & $\zh^{80}$, $u^{75}$ & modulo $2^{61}-1$ (two frames), $2^{89}-1$\\
$\cV=\varpi_0(\YCQ)_0$ (\cref{prop:con-res} with (G)) & $u^{40}$; $u^{30}$ & exact, two frames; two primes\\
$c(z;0)=\Ginf$ (\cref{thm:gam-ginf}) & $z^{60}$ (chart at $L=0$); $z^{80}$ (the value) & exact\\
$\fc^{\mathrm{can}}_m=\coef{z^{5m}}\Ginf$ (\cref{cor:gam-free}) & $m\leq 25$; $m\leq 100$ & exact; two primes\\
\end{tabular}
\caption{Finite windows of $\Dn=1$ and of the identity \eqref{eq:gam-Gamma}.}
\label{tab:evi-windows}
\end{table}

% Evidence appendix, part 3 of 3: the model. Last part of the appendix.

\subsection{The model}
\label{sec:evi-model}

\Cref{tab:evi-model} lists the main checks of \cref{sec:mod}. Further, the formula of \cref{rmk:mod-hock} was checked for $h\leq 12$ exactly, and the singular lists of the Model Lemma at genus $0$ to $3$ all vanish, while deliberately wrong frames give non-zero values. Off $\sK$ (\cref{rmk:mod-offK}), exactly at $v=1$, $\langle\psi^5H\rangle_3=53/1209600$, whereas the one-point law of \cref{cor:mod-onepoint} would give $-1/604800$.

\begin{table}[ht]
\centering\small
\begin{tabular}{>{\raggedright\arraybackslash}p{0.4\textwidth}|>{\raggedright\arraybackslash}p{0.36\textwidth}|>{\raggedright\arraybackslash}p{0.16\textwidth}}
statement & range & arithmetic\\ \hline
topological recursion on $x=\Wv$ is $(-5)^{h-1}$ times the model CohFT, i.e.\ the model correlators are $(-5)^{1-h}$ times those of topological recursion (\cref{thm:mod-tr}) & $72120$ colour-basis correlators with $\psi$-deficit $\leq 6$ at $(h,n)=(0,3)$, $(0,4)$, $(0,5)$, $(1,1)$, $(1,2)$, $(1,3)$, $(2,1)$, $(2,2)$, $(3,1)$ ($4665$ non-zero); on $x=\kappa\Wv$, $\kappa=1,4,-1$: $6495$ each & two primes\\
the R-matrix of the curve $x=\kappa\Wv$ is $\Mo(z/(\kappa\varsigma^4))$ (\cref{prop:mod-Rtr}) & all $25$ entries through $z^{10}$, $\kappa=1,4,-1$, $\varsigma=1,2$ & exact, $\bbQ(\zeta)$\\
the dilaton-leaf sign of \cite[Def.~2.13]{PPZ15} & $(h,n)=(1,1)$, four points, $\kappa=1,4$; the opposite sign fails at all four & exact\\
$\FTR_h(v=1)=B_{2h}/(2h(2h-2))=-\tfrac{1}{240},\tfrac{1}{1008},-\tfrac{1}{1440},\tfrac{1}{1056}$ & $h=2,\dots,5$ & $\bbQ(\zeta)$ ($h\leq 3$); two primes\\
the period identity of \cref{prop:mod-period}; $\Sigma_h=-4^hB_{2h}/(2h)$ & $h\leq 7$; $h\leq 20$ & exact\\
the lattice theorem (\cref{thm:mod-lattice}) & weight $\leq 400$; $\leq 1000$; $16200$ raw-leg kernel vectors & exact; two primes\\
the leg dictionary (\cref{prop:mod-LD}) & $35000$ pairings per prime; $52$ frame one-point functions, $h=2,3$ & two primes; $\bbQ(\zeta)$\\
the Model Lemma on $\sK$ (\cref{thm:mod-MLK}) & $5289$ correlators with $h\leq 4$, including $1302$ at $(h,n)=(2,4)$; $1395$ singular, all zero & exact; two primes\\
\end{tabular}
\caption{Checks on the model.}
\label{tab:evi-model}
\end{table}


